% ------------------------------------------------------------------------------
% Appendix A — drafted Phase C from Direction/outputs/withhold_opportunity/
%   stage2/3/4/4b findings; Exhibit: T4
% ------------------------------------------------------------------------------
\section{An Earlier Opportunity-Set Design}\label{app:withhold_opportunity}

Before the matched interval design of Section~\ref{sec:design}, the
project's first, registered-in-advance pass at the payment-sensitivity question used
an opportunity-set construction: flag intervals in which the unit's withholding
could plausibly cause a direction (rivals near the security boundary), match
them to comparison intervals by coarsened exact matching on system state, and
test whether a binary withheld indicator interacts with the compensation
price,
\[
\mathit{withheld}_{it} = \lambda\,(\dt \times \mathit{opp}_{it}) +
\phi\,\mathit{opp}_{it} + \kappa\,\mathit{srmc}_{it} + \alpha_i +
\eta_{q(it)} + \omega_{h(it)} + \varepsilon_{it},
\]
with unit, non-synchronous-quintile, and hour-block fixed effects, clustered
by month.

Table~\ref{tab:withhold_opp} reports the result: null. The pooled interaction
is $0.00024$ ($t=1.55$, $p=0.14$), with only TORRB4 marginal ($p=0.07$) and
the sister units on either side of zero. The June-2022 gap in the
compensation-price series, 19\% of opportunity intervals, is handled by
imputation at \$164.38/MWh as a robustness variant with the same conclusion.
A sign-flip diagnostic confirms TORRB3's wrong-signed estimate is
within-sample noise, not composition: the three Torrens units' opportunity
sets are byte-identical (the essentiality state is a station-level object),
so cross-unit differences on identical regressor sets are sampling variation.
Pelican Point's cell is uninformative --- its opportunity intervals fall in a
narrow calendar window with essentially no $\dt$ variation.

Two design features explain why this test is superseded rather than pooled
with Section~\ref{sec:results}. First, the outcome is a binary withheld
indicator, and Torrens's withheld rate inside the matched strata is already
high on both sides of the contrast (90--94\% in comparison intervals,
97--99\% in opportunity intervals): the extensive margin is nearly saturated,
so the test can only detect movement in a thin residual band. The interval design's
continuous cheap-capacity share carries the intensive margin where the
estimated response lives. Second, the opportunity flag conditions on rivals'
proximity to the boundary more coarsely than the CEM competition bins.
The null is nonetheless informative, and consistent with the mechanism
picture: on the extensive margin the posture does not flex with the prize.

\begin{table}[!htbp]\centering
\caption{Appendix: earlier opportunity-set design --- $d_t \times$ opportunity interaction (withheld indicator)}
\label{tab:withhold_opp}
\begin{adjustbox}{max width=\textwidth}
\begin{threeparttable}
\begin{tabular}{lccc}
\toprule
 & $d_t \times$ opp & $d_t^{robust} \times$ opp & $N$ \\
\midrule
Pooled & $0.00024$ & $0.00019$ & 177,854 \\
 & (0.00016) & (0.00013) & \\
TORRB2 & $0.00021$ & $0.00030$ & 57,677 \\
 & (0.00027) & (0.00021) & \\
TORRB3 & $-0.00018$ & $-0.00013$ & 57,677 \\
 & (0.00027) & (0.00024) & \\
TORRB4 & $0.00055^{*}$ & $0.00034$ & 57,677 \\
 & (0.00029) & (0.00026) & \\
PPCCGT\tnote{a} & $0.00061$ & $-0.00125$ & 4,823 \\
 & (0.00129) & (0.00400) & \\
\bottomrule
\end{tabular}
\begin{tablenotes}\footnotesize
\item Linear probability model: withheld $\sim d_t \times$ opportunity $+$ SRMC with unit, non-synchronous-quintile, and hour-block fixed effects on the CEM-matched sample, clustered by month. $d_t^{robust}$ imputes the June-2022 compensation-price gap at \$164.38/MWh. This earlier, coarser design is superseded by the specification of Section~\ref{sec:design}; it is reported for completeness.
\item[a] PPCCGT is underpowered in this design (few opportunity intervals); estimates are unstable.
\item $^{*}p<0.10$, $^{**}p<0.05$, $^{***}p<0.01$ (analytic).
\end{tablenotes}
\end{threeparttable}
\end{adjustbox}
\end{table}

